%% Object-Relative Critics for Pre-Execution Grasp Candidate Selection
%% Compiled with: pdflatex paper_risk_gated_vla.tex (or latexmk -pdf paper_risk_gated_vla.tex)
%% Converted from paper_risk_gated_vla_draft.md -- see that file for the working/annotated version.
%%
%% Required packages (all in standard TeX Live):
%%   IEEEtran, graphicx, booktabs, amsmath, amssymb,
%%   microtype, hyperref, cleveref, xcolor, multirow, array
%%
%% TARGET VENUE (decided 2026-08-02): Robotica (Cambridge University Press), non-open-access
%% (subscription) route -- chosen over IEEE T-RO/Advanced Robotics (already assigned to this
%% project's other two papers) and over stretch-tier options (RAS, Autonomous Robots, IEEE T-ASE,
%% IJRR) as the most realistic fit for this paper's actual weight (one real, twice-replicated
%% positive result + substantial honest negative/incomplete results). Confirmed SCIE-indexed,
%% confirmed hybrid with a no-cost subscription (non-OA) publishing option.
%% Per Robotica's own author instructions (cambridge.org/core/journals/robotica/information/
%% author-instructions/preparing-your-materials, checked 2026-08-02):
%%   - Abstract must be <=250 words, single paragraph -- currently 241 words (re-verify after any
%%     future abstract edit; this note goes stale easily, don't trust it blindly).
%%   - Citation style: numbered, square brackets, consecutive -- already matches IEEEtran's default,
%%     no change needed.
%%   - Required declarations added below the Conclusion: Financial Support, Competing Interests, and
%%     Data Availability Statement filled in 2026-08-04 per user decisions (no funding source;
%%     none; available from corresponding author on request, not a public repo). Author
%%     Contributions is STILL PLACEHOLDER TEXT -- needs real author names/roles, not fillable here.
%%   - LaTeX class: Robotica requires `ROB-NEW.cls` for the final typeset version, but their own
%%     instructions state the LaTeX source is NOT required for initial submission -- only a PDF is,
%%     generated from any reasonable format. The class file is requested only upon provisional
%%     acceptance. This file therefore stays on IEEEtran for now; do NOT attempt to reformat into
%%     ROB-NEW.cls before acceptance, and do not assume this PDF fails to meet initial-submission
%%     requirements just because it isn't in the journal's final class.
%%   - No explicit overall manuscript length/page limit was found in Robotica's public author
%%     instructions as of this check (two independent fetches of their author-instructions pages
%%     came back without one) -- this pass did not cut body content for a page target that may not
%%     exist; re-verify this directly with Robotica (or in the full PDF guide referenced from their
%%     site) before assuming no limit applies.

\documentclass[journal]{IEEEtran}

%% --- Core packages ---
\usepackage[T1]{fontenc}
\usepackage[utf8]{inputenc}
\usepackage{lmodern}
\usepackage{microtype}

%% --- Math ---
\usepackage{amsmath}
\usepackage{amssymb}

%% --- Graphics ---
\usepackage{graphicx}
\graphicspath{{figures/}}

%% --- Tables ---
\usepackage{booktabs}
\usepackage{multirow}
\usepackage{array}

%% --- Color ---
\usepackage{xcolor}
\definecolor{posblue}{HTML}{2166ac}
\definecolor{negred} {HTML}{d6604d}
\definecolor{neutgray}{HTML}{888888}

%% --- Hyperlinks ---
\usepackage[hidelinks]{hyperref}
\usepackage{cleveref}

%% --- Misc ---
\usepackage{balance}

% ─────────────────────────────────────────────────────────────────────────────
% Title deliberately does not lead with "counterfactual": the term has at
% least three incompatible meanings in the current literature (enumerated in
% sec:critic), and this paper's within-scene, same-episode candidate
% comparison is not the formal potential-outcomes sense a reader may assume
% from a title alone. The method keeps its full name in the body, where that
% distinction is stated before the term is relied on.
\title{Object-Relative Critics for Pre-Execution Grasp Candidate Selection}

% Double-blind: author block omitted, matching paper_tro.tex/paper_final.tex's convention
% \author{...}

% ─────────────────────────────────────────────────────────────────────────────
\begin{document}

\maketitle
\thispagestyle{empty}
\pagestyle{empty}

% =============================================================================
\begin{abstract}
We propose an \textbf{object-relative counterfactual critic} for
pre-execution grasp candidate selection: pose relative to the target
object, point-cloud statistics, and object identity, trained with a
within-scene Bradley--Terry-style pairwise loss on causally-admissible
features only. Evaluated on two independent, disjoint, held-out scene
batches, it significantly outperforms a geometric heuristic baseline on
both: $+15.6$pp ($p{=}0.00258$, $n{=}90$) on an independent
development-test batch, and $+14.0$pp ($p{=}3.24\text{e-}4$, $n{=}150$) on
a frozen confirmatory batch never inspected before evaluation -- a real,
twice-replicated effect. A complementary offline decomposition explains
where and why the critic works. We validate this result under an unusual
level of scrutiny: applying a causal-validity audit criterion from
companion work to a predecessor pipeline in the same problem family, we
found and corrected two evaluation defects that had made an earlier,
similar critic look far stronger than it actually was -- a seed-coupling
defect that broke its paired evaluation design, and a success criterion
that counted transient gripper contact regardless of whether the object
was lifted. Re-evaluated under a corrected, genuinely paired harness, that
predecessor checkpoint's apparent signal collapses to chance
(AUROC$=0.4996$); our own critic's harness is built directly on this
correction. We report negative results with equal rigor: an
uncertainty-based risk gate adds no measurable benefit, and the pairwise
loss term's independent contribution is quantifiably unresolvable at the
current data scale. The critic, a causal-validity audit tool, a
paired-evaluation harness, and a real-hardware validation protocol are
available to support future work (see Data Availability Statement).
\end{abstract}

% =============================================================================
\section{Introduction}
\label{sec:intro}

Pre-execution grasp candidate scoring -- training a model to rank a pool of
not-yet-executed grasp poses so the best one can be selected before any
physical or physically-simulated commitment -- is a common pattern across
analytic scorers~\cite{mahler2017dexnet}, learned
discriminators~\cite{graspgen2025}, and geometric rerankers, including this
project's own LGGSN (companion T-RO work). Its correctness rests on two
assumptions that are easy to state and easy to violate silently: every
input feature must be computable \emph{before} the chosen candidate
executes (causal validity), and comparing two scoring methods requires
evaluating them against the same candidate pool under the same conditions
(paired evaluation). Neither violation is visible from an aggregate accuracy
or success-rate number alone -- both look like a real result until traced
to source.

We propose an \textbf{object-relative counterfactual critic} for this
problem: pose features expressed relative to the target object,
point-cloud statistics, and object identity, every field registered
PRE\_EXECUTION-admissible \emph{before} training, not after, trained with a
within-scene Bradley--Terry-style pairwise loss on causally-admissible
features only. Evaluated on two independent, disjoint, held-out scene
batches -- one of them a frozen confirmatory batch never inspected before
its pre-registered gate was evaluated -- it significantly outperforms a
geometric heuristic baseline on both (\cref{sec:positiveresult}): a real,
twice-replicated effect, not a single lucky sample.

We hold this result to an unusual level of scrutiny. A pre-execution grasp
critic (``world model'') in the same problem family, trained on simulated
MuJoCo outcomes, was previously reported to beat a geometric heuristic
baseline by $+15.6$ percentage points. Auditing that claim before building
on it -- applying this project's own PRE\_EXECUTION-admissibility criterion
and provenance registry (companion T-RO work) to that predecessor pipeline
-- surfaced two compounding defects, neither itself a feature-provenance
violation: the evaluation harness's random seed encoded which method was
under test, so ``geometry'' and ``critic'' trials never executed against
the same scene or candidate pool despite the paired statistics assuming
they did; and the success label counted transient post-close gripper
contact, checked before any lift, regardless of whether the grasp actually
succeeded. Re-evaluated under a from-scratch, causally-admissible,
genuinely paired harness, that predecessor checkpoint's apparent signal is
AUROC$=0.4996$ against 1,500 real per-candidate outcomes -- exactly chance
(\cref{sec:diagnosis}). Our own critic's evaluation harness
(\cref{sec:harness}) is built directly on top of this correction, not
around it -- the same discipline that caught an invalid predecessor result
is what lets us trust our own.

We report our positive result with that discipline, and we report, with
equal weight, what this investigation could not establish: an
uncertainty-based risk gate with no measurable benefit, a pairwise-loss
ablation quantifiably not resolvable at the current evaluation scale
(\cref{sec:ablation}), and a small-data imitation-learning integration that
has not yet demonstrated closed-loop robustness.

This paper's evidence chain -- propose a critic, validate it under audit,
honestly characterize what doesn't work -- mirrors this project's companion
T-RO investigation into the same class of evaluation failure in a
different pipeline. We do not claim the specific defects found in the
predecessor pipeline are universal; we claim the failure mode itself -- an
aggregate metric that cannot distinguish a causally valid, fairly evaluated
result from an invalid one -- is a structural risk worth checking before
trusting any pre-execution critic's output, and that our own critic's
validation here demonstrates one concrete way to check it.

\textbf{Relation to companion T-RO work.} The PRE\_EXECUTION-admissibility
criterion and its audit tooling are this project's own prior formalization,
introduced and validated in a companion T-RO submission on a different,
retrospective pipeline (cross-embodiment LGGSN reranking). This paper's
contribution is not the criterion itself, but its \emph{prospective}
application to a second, independent pipeline, plus a new critic
architecture built causally-admissible from the start and validated with
two disjoint held-out test batches -- see \cref{sec:related} for how the
two papers are positioned to avoid overlapping claims.

% =============================================================================
\section{Related Work}
\label{sec:related}

\textbf{Learned grasp-candidate scoring and reranking.} Dex-Net~2.0
\cite{mahler2017dexnet} and GraspNet-1Billion \cite{graspnet1b2020} learn
grasp-quality predictors from large annotated/simulated datasets assuming
full scene geometry at query time; 6-DOF GraspNet
\cite{mousavian2019sixdof} extends this to a variational proposal
generator. None of the three states an input/label admissibility
distinction as a formal, checkable criterion. GraspGen \cite{graspgen2025}
is the closest prior system to satisfy the PRE\_EXECUTION-admissibility
criterion we apply here by construction: its discriminator's inference-time
inputs are a point-cloud encoding and candidate pose only, with
execution-derived signal confined to training labels. As in this project's
companion T-RO work, we treat GraspGen as a confirmed-compliant reference
case, not a counterexample -- external validation of the criterion, not
prior art for it. Pairwise ranking for grasp selection itself is not new --
a 2018-filed patent \cite{toris2021pairwisegrasp} trains a binary classifier
on pairwise comparisons derived from human-provided grasp preferences,
using hand-designed heuristic features. The distinction from our approach
is the supervision source and the criterion, not the pairwise mechanism
itself: our pairwise signal is derived from physical execution outcomes
within a scene, not human preference judgments, and every input feature is
checked against the PRE\_EXECUTION-admissibility criterion before training
-- a formal, checkable distinction absent from this and the other prior
systems surveyed here. Concurrent with this work, ``Robot Critics that
Sweat the Small Stuff'' \cite{sweatsmallstuff2026} fine-tunes a
vision-language critic on pairwise progress supervision from success/
failure rollouts and pairs it with an action-conditioned video model to
predict each candidate action's visual effect before selecting among them
-- the same pairwise-supervised candidate-selection pattern as our critic,
in a different modality (VLM/video prediction over raw visual rollouts
rather than a small MLP over point-cloud/pose features) and without a
stated input-admissibility criterion; their video-model-predicted
``visual effect of a candidate action'' is exactly the kind of quantity
our PRE\_EXECUTION criterion would need to audit case-by-case; we do not
audit it here since we did not evaluate their system.

\textbf{Evaluation-protocol failures in offline policy/critic assessment.}
Data leakage and unfair offline comparison are recognized general risks in
sequential-decision-making evaluation. Concurrent work on offline policy
evaluation for manipulation, ``Offline Policy Evaluation for Manipulation
Policies via Discounted Liveness Formulation''
\cite{offlinerl_leakage2026}, diagnoses a \emph{different} integrity
problem from ours: finite-horizon episodes bias value estimates under
classical TD/Monte Carlo evaluation (truncation bias), which they address
with a liveness-based Bellman operator, evaluated on VLA and
diffusion-policy manipulation tasks plus cloth folding. This is adjacent
evidence that offline/simulated manipulation-policy evaluation is
error-prone in more than one structurally distinct way -- not directly
overlapping prior art for the seed-coupling and success-criterion defects
\cref{sec:diagnosis} diagnoses, and should not be read as a leakage paper.
Independently, and closer in spirit to our diagnosis, ``What Are We
Actually Benchmarking in Robot Manipulation?''
\cite{jiang2026benchmarking} audits LIBERO and CALVIN and finds only
$19.8\%$ of LIBERO's reported advances are statistically significant once
evaluation noise is accounted for, alongside shortcut-solvable benchmark
items and train/test-distribution proximity effects that inflate apparent
generalization. Their specific failure modes (shortcut solvability,
creeping overfitting, data-source dependence) are distinct from the
seed-coupling and success-criterion defects \cref{sec:diagnosis} diagnoses
in our own pipeline, but the broader concern is the same and, as of
mid-2026, actively converging from multiple independent directions: a
robot-manipulation benchmark score is not self-certifying, and treating one
as evidence of general capability without auditing how it was produced is
an increasingly recognized, not a niche, risk. We position this paper as
one concrete, fully-traced instance of that broader concern, not as its
sole discovery.

\textbf{Risk-aware and uncertainty-gated action selection in VLA
policies.} ReconVLA \cite{chen2026reconvla} applies conformal prediction
directly to a pretrained VLA policy's action tokens, yielding calibrated
uncertainty estimates that correlate with execution quality, and extends
the same idea to state-space outlier detection as a failure-anticipation
mechanism. SafeVLA \cite{zhang2025safevla} frames safety alignment as a
constrained Markov decision process, optimizing against elicited unsafe
behaviors and reporting an $83.58\%$ reduction in cumulative
safety-violation cost versus prior methods. Both report a positive effect
from their respective uncertainty/safety mechanism. Our own
ensemble-disagreement risk gate -- calibrated on held-out data and
evaluated with the same statistical rigor as our positive result -- shows
no measurable benefit over the ungated critic (coverage $98.7\%$,
\cref{sec:limitations}). We do not read this as evidence against
uncertainty-gating in general: ReconVLA's per-token conformal calibration
and SafeVLA's constrained-optimization framing are both structurally
different from a post-hoc ensemble-disagreement filter applied to a
pre-execution candidate scorer. Concurrent work directly asking when
uncertainty-gating helps, ``Confidence-Gated Robot Autonomy''
\cite{gaus2026confidencegated}, identifies a dataset-dependent
\emph{competence regime}: below it, uncertainty rankings are weak and
unstable regardless of which uncertainty source (softmax, MC dropout,
ensemble) is used; above it, simple proxies suffice and the gating
threshold matters more than the method. Our critic's $\sim\!50\%$ pooled
success rate (\cref{sec:results}) may simply not yet be in the regime their
analysis identifies as necessary for a gate to add value -- we flag this as
the most likely explanation for our null result rather than leaving it
unexplained, without claiming to have tested for the regime directly. We
report our own result as a specific negative finding for the specific gate
design tested, with the same discipline as our positive finding, not
omitted or hedged into ambiguity.

\textbf{Sim-to-real transfer for grasp success prediction.} Out of scope
for this paper's evidence base; motivates \cref{sec:hardware}'s protocol as
a stated future-work step, not a claim made here.

% =============================================================================
\section{Method}
\label{sec:method}

\subsection{Causal-validity criterion (reused, not reinvented)}
\label{sec:criterion}

We reuse Definitions 1--3 and the audit algorithm from this project's
companion T-RO formalization at the level of generality needed here:
a scene induces a candidate pool, each candidate carries generator-time
attributes, and a feature is PRE\_EXECUTION-admissible only if it is
computable from information available before the candidate it describes
has executed -- with training labels themselves exempt (they are
necessarily execution-derived by definition, and only input features are
constrained). This paper's contribution here is \emph{application and
extension} (registering a new pipeline's fields, \cref{sec:critic}), not
the criterion itself.

\subsection{Object-relative counterfactual critic}
\label{sec:critic}

The candidate pose is expressed relative to the detected/simulated object
position, with sin/cos yaw, gripper opening, 9-dimensional point-cloud
statistics, and an object-identity one-hot (3 classes: CrackerBox,
MustardBottle, PowerDrill). All fields are PRE\_EXECUTION-admissible per
the \cref{sec:criterion} criterion -- a checked property, registered in
the project's feature-provenance registry, not an assumption. The
architecture is a 2-hidden-layer MLP ($64$/$64$, SiLU, dropout $0.05$),
trained as a 5-seed ensemble. The loss is binary cross-entropy on the
per-candidate ground-truth outcome, plus, for the \texttt{object\_counter\-
factual} variant, a within-scene Bradley--Terry/BPR-style pairwise term
comparing successful vs.\ failed candidates from the \emph{same} scene --
the ``counterfactual'' framing: what would have happened had a different
candidate in the same, already-observed scene been chosen. This is a
same-scene, within-episode alternative used as a \emph{training signal} for
the critic, not the post-hoc explanation or formal potential-outcomes sense
of ``counterfactual'' common in the 2025--2026 causal-ML and
explainability literature -- the pairwise comparison never requires
estimating an outcome under an unobserved intervention, only comparing two
candidates both already scored against the same, already-observed scene. A
third, distinct usage appears in concurrent work on policy robustness
\cite{cfnbc2607}, where ``counterfactual'' means holding an expert's action
fixed while perturbing visual nuisance factors -- closer to robustness
auditing than to our within-scene candidate comparison. We flag all three
usages here because the term is not converging on one meaning across this
literature, and a reader coming from any one of them should not assume
this paper means the same thing.
The training/validation split is scene-grouped (never splitting a scene's
candidates across train/val), stratified per object, and reproducible per
seed.

\subsection{Paired candidate-selection evaluation harness}
\label{sec:harness}

One candidate pool (object placement plus $K$ sampled grasp poses) is built
before any method is chosen; every compared method executes against a
\emph{fresh, identically re-placed} copy of the same scene, sharing the
same pool. This directly fixes the predecessor pipeline's seed-coupling
defect (\cref{sec:diagnosis}). The production execution primitive is the
same bilaterally-gated, weld-based physics grasp routine used elsewhere in
this project's simulation stack, not a bespoke, more lenient
reimplementation (\cref{sec:diagnosis} shows why that distinction
matters). Statistics are McNemar's exact test (paired), bootstrap
confidence intervals, and per-object and pooled effects, with all gates
pre-registered before any result was inspected.

\subsection{Why we rebuilt the evaluation (the predecessor pipeline's
defects)}
\label{sec:diagnosis}

This subsection is itself a finding, not just setup, reported with the
same discipline as a positive result rather than glossed over -- matching
this project's established house style of reporting a self-correction as
a first-class finding, not suppressing it. A predecessor pipeline's
evaluation harness had two compounding defects, neither itself a
feature-provenance violation. First, a \textbf{seed-coupling defect}: the
random-number-generator seed depended on which method was under test,
provably by construction and confirmed directly -- of 250 paired trials,
zero shared even the same executed candidate position between the
``geometry'' and ``critic'' arms, despite the paired statistics assuming
they did. Second, a \textbf{success-criterion defect}: the legacy
\texttt{contact or grasped or lifted} criterion saturated (150/150 sampled
candidates ``succeeded'' regardless of pose, across 2 objects and 15
scenes), replaced here by bilateral contact plus a verified post-settle
lift via the production grasp primitive (\cref{sec:harness}). Together,
these mean the predecessor pipeline's originally-reported effect
($+15.6$pp, ``world model'' vs.\ geometry) is unsupported by either the
pairing or the underlying success labels -- not a subtle statistical
concern, a fully saturated label and a broken pairing. This is why our own
critic's positive result (\cref{sec:positiveresult}) is built on this
from-scratch corrected pipeline rather than a patch to the original one,
and why \cref{sec:stalegate} reports the predecessor checkpoint's own
corrected, chance-level performance before presenting our own critic's
result.

% =============================================================================
\section{Results}
\label{sec:results}

\subsection{Stale-checkpoint gate (negative, reported first)}
\label{sec:stalegate}

Matching this project's own pre-registered-gate discipline, we report the
negative result before the positive one. The original critic checkpoint,
re-evaluated under the corrected pipeline, fails the pre-registered gate:
pooled effect $-14.0$pp (wrong sign vs.\ the required $\geq+8$pp),
AUROC$=0.4996$ on 1,500 real per-candidate labels (chance level). This is
not a mysterious failure: \cref{sec:diagnosis}'s success-criterion
defect contaminated its training labels too.

\subsection{Object-relative counterfactual critic -- the positive result}
\label{sec:positiveresult}

\begin{table}[t]
\centering
\caption{Paired live-executed results: geometry vs.\ the corrected,
object-relative counterfactual critic.}
\label{tab:positiveresult}
\begin{tabular}{@{}lrrrr@{}}
\toprule
Batch & $n$ & Geometry & Critic & $\Delta$ / exact McNemar \\
\midrule
Dev-test & 90 & 33.3\% & 48.9\% & $+15.6$pp, $p{=}0.00258$ \\
Confirmatory & 150 & 36.0\% & 50.0\% & $+14.0$pp, $p{=}3.24\text{e-}4$ \\
\bottomrule
\end{tabular}
\end{table}

\Cref{tab:positiveresult} summarizes both evaluations. The independent
development-test batch (base seed 200, $n{=}90$) shows geometry at
$30/90$ ($33.3\%$) vs.\ the critic at $44/90$ ($48.9\%$), $+15.6$pp
(exact McNemar $p{=}0.00258$). The frozen confirmatory batch (base seed
300, $n{=}150$, never inspected before its pre-registered gate) shows
geometry at $54/150$ ($36.0\%$) vs.\ the critic at $75/150$ ($50.0\%$),
$+14.0$pp (27 wins/6 losses, exact McNemar $p{=}3.24\text{e-}4$). Both
batches' scene keys are disjoint from training (base seed 100) and from
each other, independently verified via pairwise key-intersection checks,
not merely asserted.

\textbf{Per-object breakdown (live-executed, geometry / counterfactual).}
Development test: CrackerBox $4/30$ vs.\ $4/30$ (tie), PowerDrill $7/30$
vs.\ $18/30$, MustardBottle $19/30$ vs.\ $22/30$. Frozen confirmatory:
CrackerBox $9/50$ vs.\ $9/50$ (tie, unchanged), MustardBottle $29/50$ vs.\
$42/50$, PowerDrill $16/50$ vs.\ $24/50$. The pooled gain is concentrated
in PowerDrill and MustardBottle at both evaluation scales; CrackerBox is
flat at both -- the critic never selects a different top-1 candidate than
geometry for this object, at either scale, not a near-miss that a larger
sample might resolve.

\textbf{Live-executed vs.\ offline-re-scored reporting convention.}
\texttt{geometry} and \texttt{object\_counterfactual} were the two methods
actually live-selected during data collection, so their reported numbers
are live-executed paired outcomes -- the real online trial, not a re-run.
The BCE-only ablation (\cref{sec:ablation}) and \cref{sec:mechanistic}'s
multi-head decomposition were never live-selected in this collection run;
their numbers are offline re-scores against the same scene's fully-swept
candidate ground truth. This distinction matters for a subtler reason than
convenience: MuJoCo's contact solver is not perfectly reproducible on
marginal grasps. Independently verified across the confirmatory batch's 300
live-vs-reswept comparisons, exactly 2 flipped (a $\sim\!0.67\%$
marginal-grasp non-determinism rate; the development-test batch shows the
same pattern, 1 flip per method out of 90) -- a genuine
physical-reproducibility floor on any single-run success-rate estimate in
this simulator, not a script bug, and the same class of finding
independently documented in a companion pipeline ($\sim\!12.5\%$
single-run instability). We report it as a stated limitation rather than
treating any single execution as ground truth.

\subsection{Ablation: is the pairwise loss load-bearing?}
\label{sec:ablation}

Object-relative BCE alone (no pairwise term) is statistically
indistinguishable from the full \texttt{object\_counterfactual} variant on
the frozen confirmatory batch (2 wins/3 losses, $p{=}1.0$). This is not
merely ``$n$ is small'' -- it is quantifiably not resolvable by modestly
scaling up the same evaluation: exact power analysis on the observed
2-vs-3 discordant split gives essentially zero post-hoc power and a
required \textbf{210 discordant pairs} for $80\%$ power at this effect
size -- against only 5 observed here out of 150 evaluated pools, resolving
this would need on the order of several thousand additional evaluated
candidate pools, not a modest top-up. We do not attempt to close this
ablation by collecting more data at the current scale, and scope this
paper's central claim to avoid depending on it (\cref{sec:conclusion}): the
defensible claim is ``object-centric, causally-admissible scoring beats
geometry,'' not ``the within-scene pairwise loss term is what makes it
work.''

\subsection{Post-hoc finding: a candidate-level point-cloud feature defect,
found and fixed after \cref{tab:positiveresult}'s collection}
\label{sec:pcfix}

After \cref{tab:positiveresult}'s data was collected, inspection of the
feature-assembly code found a real representation defect, independent of
every other diagnosis in this paper: the point-cloud statistic
(\texttt{pc\_stats\_before}) is computed once per scene, before any
candidate is sampled, and was then attached identically to \emph{every}
candidate in that scene -- the critic had no candidate-specific point-cloud
geometric signal beyond pose. \Cref{tab:positiveresult}'s $+15.6$pp /
$+14.0$pp numbers were obtained with this defect still present; they show
that object-relative pose and scene context beat geometry, not that the
point-cloud feature was contributing anything beyond scene identity.

\textbf{Fix.} A new candidate-local statistic
(\texttt{pc\_stats\_local}, same 9-dim layout, cropped to a radius around
each candidate's own pose rather than the whole scene) replaces the shared
value, verified admissible under \cref{sec:criterion}'s own
PRE\_EXECUTION criterion and covered by unit tests proving: candidates with
different local geometry receive different features, identical geometry
remains deterministic, feature dimension/order are unchanged, no
execution-derived information leaks in, and the old shared-feature behavior
cannot regress unnoticed.

\textbf{Validation.} A controlled, paired replication isolates the fix as
the only variable: two \texttt{object\_counterfactual} ensembles (5 seeds,
400 epochs, unmodified training code) trained on the identical 120-scene
training population (base seed 100), differing only in whether
\texttt{pc\_stats\_local} was present or stripped before training. Evaluated
on a freshly collected 90-scene development-test batch (base seed 200,
same protocol as \cref{tab:positiveresult}'s dev-test), offline-re-scored
against real, physically executed per-candidate outcomes (the same
convention already used for \cref{sec:ablation}'s and
\cref{sec:mechanistic}'s numbers): the corrected feature reaches
$48/90$ ($53.3\%$) vs.\ the pre-fix shared feature at $36/90$ ($40.0\%$),
exact McNemar $p{=}0.0227$ (18 discordant wins for the fix vs.\ 6 against).
In this same fresh batch the corrected feature also beats live-executed
geometry ($p{=}0.0347$), while the pre-fix shared feature does not
($p{=}1.0$, statistically indistinguishable from geometry) -- consistent
with the shared feature carrying no real candidate-specific signal. The
two checkpoints' top-1 candidate choice differed in $82\%$ of scenes.

\textbf{What this does and does not license claiming.} This is real,
paired, statistically significant evidence that the fix improves ranking
quality, isolated from every confound the training data and procedure could
introduce. It is \emph{not} a replacement for \cref{tab:positiveresult}: the
fresh dev-test batch's live-executed geometry baseline ($41.1\%$) does not
match the originally reported one ($33.3\%$) at the same seed, indicating
codebase drift between the two collection dates rather than a literal
replay of the original trial, so the two tables are not directly
comparable row-for-row. The frozen confirmatory batch (base seed 300) was
deliberately not touched by this validation -- it was already read exactly
once for \cref{tab:positiveresult} and re-reading it to chase a different
number would violate this paper's own single-read discipline
(\cref{sec:positiveresult}). A confirmatory-grade claim that a
fix-incorporating critic improves on \cref{tab:positiveresult}'s deployed
numbers would require retraining on the original training population,
redeploying as the live-selected method, and evaluating against a new,
previously unused frozen batch -- not attempted here. We report the fix as
a validated correctness improvement to the representation with a measured,
significant secondary effect, not as a revision to this paper's central
claim.

\subsection{Feature ablation: is the object one-hot a shortcut?}
\label{sec:onehot}

The critic's input includes a 3-class object one-hot
(\cref{sec:critic}), which invites a fair objection: on a three-object
benchmark, an identity prior could be doing the work that the
object-relative geometry is credited with. We test this directly on the
paper's own population -- trained on the same 120-scene base-100 batch and
evaluated on the same 90-scene base-200 dev-test batch as
\cref{sec:pcfix}, offline-re-scored against real, physically executed
per-candidate outcomes, using the unmodified
\texttt{object\_counterfactual} architecture and training code (5 seeds
$\times$ 400 epochs per variant, ensemble-mean top-1 selection). The
frozen confirmatory batch was again not touched.

\begin{table}[t]
\centering
\caption{Feature ablation on the base-200 dev-test batch ($n{=}90$),
offline re-scoring. Exact two-sided McNemar against the live-executed
geometric baseline.}
\label{tab:onehot}
\begin{tabular}{lcc}
\toprule
Variant & Dev-test top-1 & vs.\ geometry \\
\midrule
Geometry (live-executed) & $37/90$ ($41.1\%$) & --- \\
Relative pose only       & $46/90$ ($51.1\%$) & $p{=}0.108$ \\
\; $+$ point cloud       & $48/90$ ($53.3\%$) & $p{=}0.0347$ \\
\; $+$ object one-hot    & $48/90$ ($53.3\%$) & $p{=}0.0433$ \\
\; $+$ both (deployed)   & $48/90$ ($53.3\%$) & $p{=}0.0347$ \\
\bottomrule
\end{tabular}
\end{table}

\textbf{The shortcut hypothesis is not supported.} Point-cloud geometry
alone and the object one-hot alone reach the same top-1 count ($48/90$
each), and neither dominates the other: they disagree on only 12 of 90
scenes, split evenly (6 each way, $p{=}1.0$). Dropping the one-hot
entirely (relative pose $+$ point cloud) retains full significance against
geometry ($p{=}0.0347$), so the reported effect does not rest on object
identity.

\textbf{What the ablation does not resolve.} At this scale the two feature
families are \emph{interchangeable rather than separable}: adding both
(the deployed variant) is indistinguishable from either alone ($p{=}1.0$
in both pairwise comparisons), so we cannot attribute the gain to one of
them. Relative pose by itself captures most of the numerical gap over
geometry ($46/90$ vs.\ $37/90$) but falls short of this paper's own
significance bar ($p{=}0.108$), and is likewise not separable from the
full variant ($p{=}0.804$). This has the same shape as the pairwise-loss
ablation in \cref{sec:ablation}: a real point estimate that the current
sample size cannot resolve into a clean per-feature attribution. We report
it as bounding what may be claimed -- the object one-hot is not carrying
the result -- rather than as a positive attribution to geometry.

\subsection{Leave-one-object-out: does point-cloud geometry generalize?}
\label{sec:loo}

\Cref{sec:onehot}'s tie between point-cloud and object-identity features
is a within-distribution result: on a closed three-object benchmark, the
two are structurally correlated, since point-cloud shape is itself
strongly predictive of which object is in the scene. We attempted to
break this collinearity directly, rather than argue about it: train on
two of the three objects only, and evaluate on the third, entirely
unseen during training -- object identity is not merely zeroed for the
held-out class here, it is undefined, so no identity shortcut is even
available. If point-cloud geometry carries real, object-general
per-candidate signal, a relative-pose-plus-point-cloud critic trained
this way should still beat pose-only on the held-out object. Same real
population as \cref{sec:onehot,sec:pcfix} (no new data collection),
\texttt{object\_counterfactual} architecture, 5 seeds $\times$ 400 epochs
per (held-out object, variant) pair, pooled across all three folds for
$n{=}90$ -- but now every prediction comes from a model that never saw
that scene's object during training. Confirmatory batch untouched.

\begin{table}[t]
\centering
\caption{Leave-one-object-out, pooled across three folds ($n{=}90$),
each scene scored by a model that never saw its object during training.}
\label{tab:loo}
\begin{tabular}{lcc}
\toprule
Variant & Held-out top-1 & vs.\ geometry \\
\midrule
Geometry (live-executed) & $37/90$ ($41.1\%$) & --- \\
Relative pose only       & $24/90$ ($26.7\%)$ & $p{=}0.0533$ \\
\; $+$ point cloud       & $27/90$ ($30.0\%)$ & $p{=}0.1102$ \\
\bottomrule
\end{tabular}
\\[0.3em]
{\footnotesize Both variants trend against geometry, not toward it; pose-only vs.\ pose$+$point-cloud: $p{=}0.7283$.}
\end{table}

\textbf{This does not resolve \cref{sec:onehot}'s tie in the hoped
direction -- it reveals a more fundamental limit.} Both leave-one-out
variants substantially \emph{underperform} geometry on the held-out
object, the opposite of the in-distribution pattern, and remain
statistically indistinguishable from each other ($p{=}0.7283$, essentially
unchanged from \cref{sec:onehot}'s in-distribution $p{=}1.0$). We cannot
conclude point-cloud geometry carries object-general signal from this
test, because neither feature variant generalizes to an unseen object at
all in this setting: the critic's in-distribution gain over geometry
appears to rest substantially on something that does not transfer, not
on object-general geometric reasoning. This independently replicates,
via a structurally different mechanism (leave-one-out among the three
training objects, not a genuinely novel category), the same qualitative
finding as this project's companion T-RO submission's zero-shot pear
probe on the same critic family -- every variant loses to geometry
out-of-distribution there too. We report this as a real negative result
bounding \cref{sec:onehot}'s claim, not as a resolution of it: the object
one-hot is not a shortcut relative to point-cloud geometry, but neither
feature demonstrably generalizes beyond the three-object training
distribution this paper evaluates within.

\subsection{Mechanistic analysis: multi-head decomposition (offline,
complementary)}
\label{sec:mechanistic}

\Cref{sec:positiveresult} establishes \emph{that} the corrected critic
beats geometry, live-executed -- the paper's primary quantitative claim.
This subsection asks \emph{why and where} it works, via a second,
explicitly \textbf{offline re-scoring} experiment on the same
causally-admissible features. It is complementary evidence, not a
replication, and must never be pooled with or presented as confirming
\cref{sec:positiveresult}'s live-executed numbers.

The same object-relative feature set is decomposed into four prediction
heads instead of one: \texttt{bilateral\_contact}, \texttt{lifted},
\texttt{success} (binary), and a 3-class \texttt{failure\_type}
(\texttt{success}/\texttt{no\_contact}/\texttt{weld\_no\_lift}). Two
additional classes originally specified, \texttt{contact\_no\_weld} and
\texttt{lifted\_then\_dropped}, were confirmed \textbf{structurally
absent} from every collected batch -- the weld gate is currently identical
to bilateral contact by construction, and the \texttt{fell\_off} threshold
never triggers within the post-grasp settle window -- so the 3-class
scheme is what was actually trained; the 6-class taxonomy is documented
future schema, not a current claim. Training reuses the same
base-100/dev-200/confirmatory-300 split chain as \cref{sec:positiveresult},
with checkpoint and loss-weighting choice frozen on dev-200 only and
confirmatory-300 read exactly once.

The success head alone, used for top-1 selection, beats geometry pooled
$+10.7$pp ($46.0\%$ vs.\ $35.3\%$, McNemar $p{=}0.0052$) on the offline
confirmatory-300 batch -- directionally and in rough magnitude consistent
with \cref{sec:positiveresult}'s live-executed $+14.0$pp, reported as
separate, complementary evidence for the same claim, not pooled or
averaged with it. All three heads reach AUROC $0.90$--$0.93$ with
ECE$\approx 0.04$ (reasonably calibrated).

\textbf{Interpretability payoff.} The failure-type confusion matrix exposes
a specific, actionable asymmetry a single-head success critic cannot show:
$92.2\%$ recall on true \texttt{no\_contact}, but $32\%$ of true successes
are misclassified as \texttt{no\_contact} -- the model is
conservative/under-confident on success, not dangerously over-confident on
failure.

\textbf{Limitation, load-bearing, not a footnote.} \texttt{weld\_no\_lift}
support is $100\%$ drill-attributed ($118/118$); with PowerDrill excluded,
training support for that class collapses to 1 example. This must be
reported as a within-drill measurement, not a demonstrated cross-object
failure-mode generalization capability. A genuine leave-one-object-out test
was not run in this pass.

\textbf{A methodological note worth keeping, not hiding.} The
\texttt{equal} vs.\ \texttt{success\_weighted} loss-weighting ablation
produced measurably different trained weights (confirmed by direct
parameter comparison) but \emph{identical} argmax-based top-1 accuracy at
every one of 5 seeds on the small internal validation split used during
training -- only the larger dev-200 batch's continuous AUROC distinguished
them. Coarse top-1 metrics can be an insufficiently sensitive selection
criterion for this kind of ablation at small validation-set scale; this was
caught by cross-checking, not silently reported as ``no difference.''

Neither this subsection nor \cref{sec:positiveresult} licenses a risk-gate
or VLA-policy improvement claim (\cref{sec:limitations}'s negative results
on both stand unchanged) -- the AUROC/interpretability gains here are
about the critic's own outputs, not about downstream gating or
imitation-policy performance.

% =============================================================================
\section{Limitations}
\label{sec:limitations}

Per this study's own rule -- do not report only positive results;
distinguish evidence from hypothesis from refuted conclusion -- we list
every negative or unresolved finding this investigation produced,
condensed to the subset that bears directly on this paper's central claim.

\begin{enumerate}
\item \textbf{Risk gate.} No measurable benefit over the ungated critic
(coverage $98.7\%$ -- the gate rarely actually fires). We do not claim a
gating contribution. The most likely explanation, per concurrent work on
when uncertainty-gating helps at all \cite{gaus2026confidencegated}
(\cref{sec:related}), is that our critic's $\sim\!50\%$ pooled success rate
is not yet in the competence regime that analysis identifies as necessary
-- stated as the most likely explanation, not a claim we separately tested
for it.
\item \textbf{Pairwise-loss contribution} unresolved, and not cheaply
resolvable at this evaluation scale (\cref{sec:ablation} -- quantified:
$\sim\!210$ discordant pairs needed for $80\%$ power, against 5 observed).
The paper's central claim does not depend on resolving this.
\item \textbf{Simulation-only}: no real-hardware validation in this paper
(\cref{sec:hardware} states the plan).
\item \textbf{Small object set} (3 YCB objects); generalization to a
broader object distribution untested.
\item \textbf{Physics-level non-determinism} on marginal grasps
($\sim\!0.6$--$1\%$ outcome-flip rate on repeated identical execution) -- a
measurement-noise floor on any single-run success-rate estimate in this
simulator, disclosed rather than hidden.
\item \textbf{Imitation-learning integration} (ACT) is a working pipeline,
not a validated capability -- first online rollout failed; 5
demonstrations/object is not enough data to expect closed-loop robustness.
Reported as integration status, not as a result bearing on the paper's
central claim.
\end{enumerate}

% =============================================================================
\section{Toward Real-Hardware Validation}
\label{sec:hardware}

We designed, but did not execute, a protocol to test whether
\cref{sec:positiveresult}'s simulated advantage transfers to physical
SO-ARM101 hardware -- a sim-to-real gap check, not a re-run of
\cref{sec:positiveresult}'s statistical claims; a physical pilot that fails
to replicate the simulated effect would not invalidate
\cref{sec:positiveresult}, and is the same class of outcome this project's
own prior real-hardware work has already hit and reported honestly.

\textbf{Object choice, corrected before any data collection.} The critic's
object one-hot only covers \{CrackerBox, MustardBottle, PowerDrill\}
(\cref{sec:critic}); an earlier proposal to pilot on Pear was caught,
during protocol design, as an out-of-distribution input the critic was
never trained to handle -- not a like-for-like generalization test. The
protocol instead specifies CrackerBox and MustardBottle, both
in-distribution and both objects this project already handles on real
SO-ARM101 hardware. This correction is recorded here, before any trial, per
this study's own no-post-hoc-selection discipline (\cref{sec:harness}).

\textbf{Platform.} SO-ARM101 with a relative-delta-clamped real backend, a
RealSense D435i for pose estimation, and the same rotated-mount geometry
and top-down IK bias already used in this study's simulation. Connectivity,
calibration, and low-level motor control are already verified working from
this project's independent real-hardware track; a thin real-hardware
analogue of the simulation's candidate-pool-build/score/execute pipeline is
new code to write, not new infrastructure to design.

\textbf{Design.} A matched-block paired trial (or, if placement restoration
proves unreliable, an explicitly-declared unpaired two-proportion design
instead -- chosen before running, not after seeing results) with a small
pre-registered pilot (12--20 total grasp attempts across both objects and
methods) gating a larger scale-up (20--30 trials/object/condition): proceed
only if the pilot's direction is consistent with simulation on both objects
and no safety incident occurred; otherwise stop and report the sim-to-real
gap honestly rather than scaling up in search of significance. A
seven-item safety checklist (continuous human supervision, an active
relative-motion clamp, pre-execution joint-range verification, workspace
clearance, a confirmed e-stop path, reduced speed through the pilot stage,
and an explicit collision check against the camera mount and tray) must be
re-confirmed live at execution time, not assumed from this document.

\textbf{What this would, and would not, settle.} A completed pilot/scale-up
would establish whether the simulated $+14$--$16$pp advantage survives the
sim-to-real gap on CrackerBox and MustardBottle specifically, at the
tested sample size. It would not establish generalization to PowerDrill or
Pear, and does not test the risk gate at all (\cref{sec:limitations}) --
only top-1 critic selection, to keep the physical pilot's scope minimal and
interpretable.

\textbf{Status, stated plainly.} Camera repositioning is an unresolved
physical-setup blocker, shared with an unrelated, independent line of work
in this project, and is a precondition for every later stage. No
real-hardware result from this section exists in this paper.

% =============================================================================
\section{Conclusion}
\label{sec:conclusion}

We proposed an object-relative counterfactual critic for pre-execution
grasp candidate selection and found a real, replicated, live-executed
effect: it beats a geometric baseline by $+15.6$pp on an independent
development batch ($p{=}0.00258$, $n{=}90$) and $+14.0$pp on a frozen,
pre-registered confirmatory batch never inspected before evaluation
($p{=}3.24\text{e-}4$, $n{=}150$) -- two independent confirmations, not one
lucky sample. We validated this result to an unusual standard: applying
and extending this project's causal-validity criterion to a predecessor
pipeline in the same problem family, we found that its own strongest
reported result was invalid -- not through a single bug, but two
independent evaluation defects that each, alone, would have been enough to
invalidate the comparison. Rather than treat that as a dead end, we
rebuilt the evaluation as a genuinely paired, causally-admissible design
and built our own critic's validation directly on top of the correction,
with a complementary offline decomposition explaining where and why the
effect concentrates. We report our positive result as the paper's central
claim, and only this claim -- not a claim about the specific pairwise loss
term, which exact power analysis shows is not resolvable at the current
data scale without roughly two orders of magnitude more evaluated
candidate pools (\cref{sec:ablation}), and not a claim about risk-aware
gating, which measurably added nothing here (\cref{sec:limitations}). A
real-hardware validation protocol is designed and the platform's
connectivity and low-level control are independently verified, but no
hardware result exists in this paper (\cref{sec:hardware} states the plan,
not a result).

The discipline we believe generalizes beyond this specific pipeline is not
the critic architecture or the loss function -- it is the practice of
auditing an evaluation's construction before trusting its output, reporting
a negative gate result with the same rigor as a positive one, and stating
plainly, at every step, what the evidence does and does not support.

% =============================================================================
% Robotica (Cambridge University Press) requires these declarations at
% submission -- see https://www.cambridge.org/core/journals/robotica/information/author-instructions/preparing-your-materials.
% PLACEHOLDER TEXT: fill in with real answers before submission; do not
% submit with placeholder content still in place.

\section*{Author Contributions}
[FILL IN: conception, methodology, software, analysis, writing -- per
author, per CRediT taxonomy or equivalent.]

\section*{Financial Support}
This research received no specific grant from any funding agency,
commercial, or not-for-profit sectors.

\section*{Competing Interests}
The authors declare none.

\section*{Data Availability Statement}
The materials supporting this study's findings --- the scene-level
evaluation batches (training, development-test, and the frozen
confirmatory batch), each recording every sampled candidate together with
its physically executed grasp outcome; the trained critic checkpoint
ensembles for every reported variant, including those of the feature and
loss ablations; and the paired-evaluation harness and analysis code ---
are available from the corresponding author upon reasonable request.

% =============================================================================
\begin{thebibliography}{10}

\bibitem{mahler2017dexnet}
J.~Mahler, J.~Liang, S.~Niyaz, M.~Laskey, R.~Doan, X.~Liu, J.~A.~Ojea,
and K.~Goldberg,
``Dex-Net 2.0: Deep Learning to Plan Robust Grasps with Synthetic
  Point Clouds and Analytic Grasp Metrics,''
\emph{Proc.\ Robotics: Sci.\ Syst.\ (RSS)}, 2017.

\bibitem{graspnet1b2020}
H.-S.~Fang, C.~Wang, M.~Gou, and C.~Lu,
``GraspNet-1Billion: A Large-Scale Benchmark for General Object
  Grasping,''
\emph{Proc.\ IEEE/CVF Conf.\ Computer Vision and Pattern Recognition
  (CVPR)}, 2020.

\bibitem{mousavian2019sixdof}
A.~Mousavian, C.~Eppner, and D.~Fox,
``6-DOF GraspNet: Variational Grasp Generation for Object
  Manipulation,''
\emph{Proc.\ IEEE/CVF Int.\ Conf.\ Computer Vision (ICCV)}, 2019.

\bibitem{graspgen2025}
A.~Murali, B.~Sundaralingam, Y.-W.~Chao, W.~Yuan, J.~Yamada,
M.~Carlson, F.~Ramos, S.~Birchfield, D.~Fox, and C.~Eppner,
``GraspGen: A Diffusion-based Framework for 6-DOF Grasping with
  On-Generator Training,''
\emph{arXiv:2507.13097}, 2025.

\bibitem{offlinerl_leakage2026}
H.~Wang, J.~Bowden, C.~Crosby, and S.~Bansal,
``Offline Policy Evaluation for Manipulation Policies via Discounted
  Liveness Formulation,''
\emph{arXiv:2605.11479}, 2026.

\bibitem{jiang2026benchmarking}
T.~Jiang, X.~Tan, S.~Wheeler, L.~Sun, T.~W.~Ayalew, and M.~Walter,
``What Are We Actually Benchmarking in Robot Manipulation?,''
\emph{arXiv:2606.04233}, 2026.

\bibitem{chen2026reconvla}
L.~Chen, Z.~Lyu, and W.~J.~Beksi,
``ReconVLA: An Uncertainty-Guided and Failure-Aware
  Vision-Language-Action Framework for Robotic Control,''
\emph{arXiv:2604.16677}, 2026.

\bibitem{zhang2025safevla}
B.~Zhang, Y.~Zhang, J.~Ji, Y.~Lei, Y.~Cai, J.~Dai, Y.~Chen, and
Y.~Yang,
``SafeVLA: Towards Safety Alignment of Vision-Language-Action Model
  via Constrained Learning,''
\emph{Proc.\ Neural Information Processing Systems (NeurIPS)}, 2025.

\bibitem{gaus2026confidencegated}
J.~A.~Gaus, J.~P.~F.~Charaja, and D.~Haeufle,
``Confidence-Gated Robot Autonomy: When Does Uncertainty Actually
  Help?,''
\emph{ICRA 2026 Workshop, arXiv:2605.18045}, 2026.

\bibitem{toris2021pairwisegrasp}
R.~Toris, S.~Elliott, and D.~Kent,
``Method and System for Selecting a Preferred Robotic Grasp of an
  Object-of-Interest Using Pairwise Ranking,''
U.S. Patent 10\,899\,011 B2, Fetch Robotics, Inc., filed 2018, issued 2021.

\bibitem{cfnbc2607}
G.~D'urso, K.~Roy, N.~Lawrance, and B.~Tidd,
``It's Not Just More Demos: Counterfactual Action Sensitivity Coverage
  for Data-Efficient Robust Robot Imitation,''
\emph{arXiv:2607.27261}, 2026.

\bibitem{sweatsmallstuff2026}
S.~Sudhakar, J.~Liang, S.~Rammohan, P.~Tokmakov, R.~Zemel, and C.~Vondrick,
``Robot Critics that Sweat the Small Stuff,''
\emph{arXiv:2606.21572}, 2026.

\end{thebibliography}

\end{document}
